% !PRINTED 476
\weberpage{476}

Als einfaches Beispiel möge hier nur noch die \gesperrt{Gruppe der cyklischen
Permutationen, kurz die cyklische Gruppe}, angeführt werden, die man
erhält, wenn man in irgend einer Anordnung von Ziffern $(1, 2, 3
\ldots n)$ jede Ziffer durch die folgende und die letzte durch die
erste ersetzt, und wenn man diese Permutation so oft wiederholt, bis
man zu der ursprünglichen Anordnung zurückkommt. Der Grad der
cyklischen Gruppe ist gleich der Anzahl der Ziffern, z.\,B. bei drei
Ziffern:
\[
\begin{pmatrix} 1, & 2, & 3 \\ 2, & 3, & 1 \end{pmatrix}, \quad
\begin{pmatrix} 1, & 2, & 3 \\ 3, & 1, & 2 \end{pmatrix}, \quad
\begin{pmatrix} 1, & 2, & 3 \\ 1, & 2, & 3 \end{pmatrix}.
\]

Es giebt besondere Gruppen, die uns später noch mehr beschäftigen
werden, die nur aus solchen Permutationen bestehen, bei deren
Zusammensetzung ausser den oben erwähnten auch noch das commutative
Gesetz gilt, bei denen also immer $\pi_1 \pi_2 = \pi_2 \pi_1$ ist, und
in Folge dessen die Analogie mit der Multiplication eine viel
vollständigere ist. Solche Gruppen, zu denen die cyklischen Gruppen
gehören, heissen \gesperrt{commutative} oder
\gesperrt{Abel'sche Gruppen}. Wenn die sämmtlichen Elemente einer
Permutationsgruppe $Q_1$ in einer Gruppe $Q$ enthalten sind, so heisst
$Q_1$ ein \gesperrt{Theiler von $Q$, und $Q$ durch $Q_1$ theilbar}.

\begin{center}
\S.\,149.

\gesperrt{Galois'sche Gruppe.}
\end{center}

Wenn wir in den Wurzeln $\alpha, \alpha_1, \alpha_2 \ldots
\alpha_{m-1}$ einer Gleichung $F(x) = 0$, wie sie durch \S.\,148, (3)
als Grössen des Normalkörpers $\Omega(\varrho)$ vom $\mu^{\mathrm{ten}}$
Grade ausgedrückt sind, eine der $\mu$ Substitutionen $\sigma$ des
Normalkörpers ausführen, so vollzieht sich unter diesen Wurzeln eine
gewisse Permutation, wie schon oben gezeigt ist. Es entspricht also
jeder Substitution $\sigma$ eine gewisse Permutation $\pi$ der $m$
Ziffern $0, 1, 2 \ldots m-1$. \gesperrt{Zwei verschiedene Substitutionen
$\sigma_h, \sigma_k$ führen auch zu zwei verschiedenen Permutationen
$\pi_h, \pi_k$.}

Denn nach \S.\,145 ist $\varrho$ eine rationale Function von $\alpha,
\alpha_1 \ldots \alpha_{m-1}$, die wir durch
\[
\varrho = \Phi(\alpha, \alpha_1 \ldots \alpha_{m-1})
\]
bezeichnen wollen. Wenn wir hierin
\[
\alpha_s = \chi_s(\varrho)
\]
setzen, so folgt

\weq{1}{\varrho = \Phi[\chi(\varrho), \chi_1(\varrho) \ldots \chi_{m-1}(\varrho)],}
%CONT>>
